\documentclass[12pt,a4paper]{article}
\usepackage[utf8]{inputenc}
\usepackage[T1]{fontenc}
\usepackage{lmodern}
\usepackage[french]{babel}

\usepackage[a4paper,vmargin={3cm,3cm},hmargin={2.5cm,2.5cm}]{geometry}
\usepackage{fancyhdr}

\usepackage{mathtools}
\usepackage{amssymb}
\usepackage{stmaryrd}
\usepackage{xcolor}

\definecolor{myblue}{RGB}{20,60,160}

\setlength{\parindent}{0pt}
\setlength{\parskip}{.45em}
\pagestyle{fancy}\fancyhf{}\fancyhead[L]{Probabilités et statistiques - L2 S1}\fancyhead[R]{M. Barry, V.Dan}\fancyfoot[C]{\thepage}

\newtheorem{exercice}{Exercice}

\newcommand{\bemph}[1]{\textbf{\textcolor{myblue}{\emph{#1}}}}
\newcommand{\intervalle}[2]{\llbracket #1,#2 \rrbracket}

\renewcommand{\P}{\mathbb{P}}
\renewcommand{\leq}{\leqslant}
\renewcommand{\geq}{\geqslant}

\newenvironment{corrige}{\par\smallskip\noindent\textbf{Corrigé.} }{\par\medskip}

\begin{document}

\begin{center}{\Large\bfseries Feuille d'exercices n\textsuperscript{o} 2 \\[.2em]{\large Probabilités conditionnelles et indépendances }}\end{center}

\begin{exercice} \textbf{( Probabilité uniforme )} \\
On jette une pièce de monnaie équilibrée trois fois de suite.
\begin{enumerate}
    \item  Décrire l'espace des événements. 
    \item   Écrire les sous-ensembles de ; correspondants aux évènements suivants : 
    \begin{enumerate}
        \item $A_i$ : " la pièce sort pile au iéme jet", 
        \item $B_i$ : " la pièce sort au plus i piles ",
        \item $C_i$ : " pile sort exactement i fois ".
    \end{enumerate}
    \item  Calculer $\mathbb{P}(A_i)$, $\mathbb{P}(B_i)$, $\mathbb{P}(C_i)$ et $\mathbb{P}(\Bar{B_i})$,  $i=1 , 2, 3$. 
\end{enumerate}
    
\end{exercice}
\begin{exercice} \textbf{( Combinatoire )} \\
On lance trois dés non truqués. Soient les événements suivants : \\
A : " On obtient au moins un 1 ", \\
B : " On obtient au moins deux faces identiques ", \\
C : " La somme des nombres obtenus est paire".
\begin{enumerate}
    \item  Calculer $\mathbb{P}(A)$, $\mathbb{P}(B)$ et $\mathbb{P}(C)$. 
    \item  Calculer $\mathbb{P}(A \cap B)$ , $\mathbb{P}(A \cap C)$ et $\mathbb{P}(B \cap C)$
\end{enumerate}
    
\end{exercice}

\begin{exercice} ( \textbf{Quelques propriétés des probabilités conditionnelles.})  \par 
Soit A et B deux événements d'un espace probabilisé $(\Omega, \mathcal{F}, \mathbb{P})$. \par
\textbf{Question 1.} Démontrer les propriétés suivantes:
\begin{enumerate}
    \item  $\mathbb{P}(A | \Omega ) = \mathbb{P}(A)$
    \item $\mathbb{P}(A \cap B | B) = \mathbb{P}(A | B)$ 
    \item Si B n'est pas négligeable et $B \subseteq A$, alors $\mathbb{P}(A | B) = 1$ 
\end{enumerate} 
\textbf{Question 2.} Soit $C \in \mathcal{F}$ tel que $\mathbb{P}( C) > 0$. 
\begin{enumerate}
    \item Montrer si $A \subseteq B$, alors $\mathbb{P}(A | C) \leq \mathbb{P}(B | C)$; 
    \item Additivité forte : Montrer que  $\mathbb{P}( A\cup B | C)= \mathbb{P}( A| C) + \mathbb{P}( B | C) - \mathbb{P}(A\cap B | C)$
\end{enumerate}
\end{exercice}

\begin{exercice} \textbf{(Probabilité totale et formule de Bayes )}

On dispose de deux urnes $U_1$ et $U_2$. L'urne $U_1$ contient trois boules blanches et une boule noire. L'urne $U_2$ contient une boule blanche et deux boules noires. On lance un dé non truqué. Si le dé donne un nombre inférieur ou égale 2, on tire une boule dans $U_1$. Sinon on tire une boule dans $U_2$. On suppose que les boules sont indiscernables au toucher. \\ 

\textbf{Question 1.} Calculer la probabilité de tirer une boule blanche. \\
\textbf{Question 2.} On a tire une boule blanche. Calculer la probabilité qu'elle provienne de l'urne $U_1$.
\end{exercice}

\begin{exercice} (\textbf{Vrai ou Faux } ) \par
Soit A, B, C trois événements d'un espace probabilisé $(\Omega, \mathcal{F}, \mathbb{P})$. Répondre par vrai ou faux aux questions, en justifiant la réponse. \\
\textbf{Question 1.} Si A et B sont indépendants, et si B et C sont indépendants, alors A et C sont indépendants. \\
\textbf{Question 2.} Si A et B ne sont pas indépendants, et si B et C ne sont pas indépendants, alors A et C ne sont pas indépendants. \\
\textbf{Question 3.} Si $A\cap B$ et $A \cap C$ sont indépendants, alors B et C sont
indépendants. \\
\textbf{Question 4.} On a : $\mathbb{P}(C | A\cup B) = \mathbb{P}(C| A) + \mathbb{P}(C | B)$ - $\mathbb{P}(C | A\cap B)$
\end{exercice}

\begin{exercice} ($\star \star \star$) \\
Soit $(\Omega, \mathcal{F}, \mathbb{P})$ un espace probabilité; $(A_n)_{n \geq 1}$ une suite décroissante d'événements sur $\Omega$ vérifiant $\mathbb{P}(A_1) = \frac{1}{2}$, et $\forall n\geq 1$, $\mathbb{P}(\Bar{A}_{n+1} | A_n) \leq \frac{1}{(n+1)^2}$. \\

\textbf{Question 1.} Montrer que : \\
\[
\forall n\geq 1, \quad \mathbb{P}(A_n) \geq \frac{1}{2}\prod_{i=1}^{n-1}( 1 - \frac{1}{(i+1)^2})
\] 
\textbf{Question 2.} Montrer que pour tout $x \in [0, \frac{1}{2} ]$, $-\log (1 - x ) \leq 2x$.\\  
En déduire que  $\forall n\geq 1,\quad \mathbb{P}(A_n) \geq \frac{1}{2e^2}$. \\

\textbf{Question 3} Montrer que $\mathbb{P}(\cap_{i=1}^{+\infty} A_i ) > 0$
    
\end{exercice}

\end{document}